% =====================================================================
% CHAPTER 2 — LITERATURE REVIEW (~1,500–1,800 words; 15–18% of the
% supplied ~10,000-word planning total)
%
% Job: assemble the FOUR shelves the layers stand on, and show the gap
% at their intersection. Organise by theme; end each section with one
% sentence of synthesis (agreed / contested / missing).
%
%   2.1 Text sentiment and asset prices   -> grounds L1
%   2.2 Sentiment extraction: lexicons -> FinBERT -> LLMs
%   2.3 LLM measurement problems          -> grounds L2/L3
%       (label flattening, calibration, look-ahead bias)
%   2.4 Differences-of-opinion theory     -> grounds L2
%   2.5 Generalizability theory           -> grounds L3
%   2.6 Synthesis and gap
%
% KEY POSITIONING MOVE (from the research plan): the multi-agent
% trading literature treats model consensus as CONFIDENCE; this
% dissertation argues consensus proxies CROWDING. Make that opposition
% explicit in 2.3/2.4 — it is the sharpest statement of novelty.
%
% NOTE: ~48 paper stubs are queued in the wiki under tag `bts-reading`
% — work through the P0 list first (Hong & Stein, Tetlock, Malo,
% Harris & Raviv, Shavelson & Webb ch.1–2, MacKinlay, Glasserman & Lin).
% =====================================================================
\chapter{Literature Review}
\label{ch:literature}

\section{Text Sentiment and Asset Prices}
\label{sec:lit-sentiment-prices}

[The canonical starting point is \textcite{tetlock2007}: media
pessimism predicts price pressure and reversal --- note that reversal
is already present in the very first paper, which foreshadows L2.
Then the modern strand: \textcite{mitic2026llm} shows an
\gls{llm}-augmented pipeline preserving 83\% autocorrelation / 93\%
volatility consistency. Contrast with early null results on
polarity-only sentiment --- the scalar being too thin from the other
side.]

\section{Sentiment Extraction: Lexicons to LLMs}
\label{sec:lit-extraction}

[\gls{vader} \parencite{hutto2014vader} and dictionary methods;
FinBERT \parencite{araci2019finbert}; generative \glspl{llm}.
Introduce the labeled-set lineage here too: PhraseBank
\parencite{malo2014phrasebank} with its 50/66/75/100\% annotator
agreement tiers --- those tiers are not a footnote, they are the
human-disagreement ground truth that validates the L2 ambiguity proxy
and the L3 tier check.]

\section{LLM Measurement Problems}
\label{sec:lit-llm-measurement}

% This section carries the methodological stakes of the dissertation.
Three strands matter. \emph{Label flattening:} \glspl{llm} compress
human label variation toward consensus answers [cite the
label-variation papers from the reading list] --- the erased
disagreement is precisely what L2 trades on. \emph{Calibration:}
modern networks are systematically overconfident
\parencite{guo2017calibration}; soft-label scoring with Brier/\gls{ece}
evaluation is the repo's response. \emph{Look-ahead bias:} models
scoring pre-cutoff text may have memorized outcomes
\parencite{glasserman2023lookahead, gao2025lookahead} --- this
motivates the leakage protocol in \cref{sec:data-leakage}.

\section{Differences of Opinion and Crowded Trades}
\label{sec:lit-disagreement}

[The theory shelf for L2: disagreement, not mean belief, drives volume
\parencite{harris1993differences, kandel1995differential}; disagreement
models of the market \parencite{hong2007disagreement}. Then the
opposition: multi-agent \gls{llm} trading systems use consensus as a
confidence signal, whereas if the industry converges on the same few
foundation models, consensus marks a \emph{crowded interpretation} ---
the same trade on the same reading. State this contrast explicitly.]

\section{Generalizability Theory}
\label{sec:lit-gtheory}

[The measurement shelf for L3: \textcite{cronbach1972} and the primer
treatment of \textcite{shavelson1991}. G-theory decomposes score
variance into text, model, prompt, sample, and interaction facets,
yielding per-item reliability --- something classification accuracy
cannot express. Recent G-theory applications to \gls{llm} outputs are
method imports, not competitors; the finance application is open.]

\section{Synthesis and Gap}
\label{sec:lit-gap}

% One paragraph, stated plainly enough that Chapter 3 reads as the
% obvious response.
[No prior work treats the ensemble of machine readings as a
measurement instrument: the sentiment--trading literature trades the
mean; the measurement literature never reaches a P/L statement. The
gap is the intersection, and RQ2/RQ3 occupy it.]
